\SourceLeaf{159}
Thus, in taking $\alpha A^2+\beta B^2+\gamma C^2+\cdots$ as the sum, one neglects all of $(a-A)^2+(a'-A)^2+\cdots$, $(b-B)^2+(b'-B)^2+\cdots$, $\cdots$, which are all positive.

The proof is easy. We have
\[
\begin{aligned}
aa &= [A+(a-A)]^2 = AA+2A(a-A)+(a-A)^2\\
a'a' &= [A+(a'-A)]^2 = AA+2A(a'-A)+(a'-A)^2\\
a''a'' &= [A+(a''-A)]^2 = AA+2A(a''-A)+(a''-A)^2
\end{aligned}
\]
\begin{center}
. . . . . . . . . . . . . . . . .
\end{center}
\[
\Sigma aa = \alpha AA+2A\Sigma(a-A)+\Sigma(a-A)^2
\]
or, since $\Sigma(a-A)=\Sigma a-\alpha A=0,$
\[
\Sigma aa = \alpha AA+\Sigma(a-A)^2.
\]
\Needspace{6\baselineskip}
\begin{center}
\textsc{Gauss} to \textsc{Olbers}. Göttingen, 3 May 1827.
\end{center}

\begin{center}
. . . . . . . . . . . . . . . . .
\end{center}

For a successful application of probability theory to observations, a thorough knowledge of the \textls{subject} is always of the highest importance. Where it is lacking, exclusion on account of a larger discrepancy is always problematic unless the number of observations available is very large. All the individual components of observational error that we cannot avoid surely have limits, even if we cannot specify them precisely. In very many cases we can say with certainty that a large error that has occurred lies outside the limits of possibility of such errors, and that an extraordinary mistake must have been made. Such cases must of course be excluded. But as long as it is conceivable that the error arose from an unfortunate combination of the components, it should not be excluded. There may occasionally be cases in which one doubts whether to count them in the first or second class. One may then proceed as one wishes,
